% Options for packages loaded elsewhere
% Options for packages loaded elsewhere
\PassOptionsToPackage{unicode}{hyperref}
\PassOptionsToPackage{hyphens}{url}
\PassOptionsToPackage{dvipsnames,svgnames,x11names}{xcolor}
%
\documentclass[
  sigconf,
  review]{acmart}
\usepackage{xcolor}
\usepackage{amsmath,amssymb}
\setcounter{secnumdepth}{5}
\usepackage{iftex}
\ifPDFTeX
  \usepackage[T1]{fontenc}
  \usepackage[utf8]{inputenc}
  \usepackage{textcomp} % provide euro and other symbols
\else % if luatex or xetex
  \usepackage{unicode-math} % this also loads fontspec
  \defaultfontfeatures{Scale=MatchLowercase}
  \defaultfontfeatures[\rmfamily]{Ligatures=TeX,Scale=1}
\fi
\usepackage{lmodern}
\ifPDFTeX\else
  % xetex/luatex font selection
\fi
% Use upquote if available, for straight quotes in verbatim environments
\IfFileExists{upquote.sty}{\usepackage{upquote}}{}
\IfFileExists{microtype.sty}{% use microtype if available
  \usepackage[]{microtype}
  \UseMicrotypeSet[protrusion]{basicmath} % disable protrusion for tt fonts
}{}
\makeatletter
\@ifundefined{KOMAClassName}{% if non-KOMA class
  \IfFileExists{parskip.sty}{%
    \usepackage{parskip}
  }{% else
    \setlength{\parindent}{0pt}
    \setlength{\parskip}{6pt plus 2pt minus 1pt}}
}{% if KOMA class
  \KOMAoptions{parskip=half}}
\makeatother
% Make \paragraph and \subparagraph free-standing
\makeatletter
\ifx\paragraph\undefined\else
  \let\oldparagraph\paragraph
  \renewcommand{\paragraph}{
    \@ifstar
      \xxxParagraphStar
      \xxxParagraphNoStar
  }
  \newcommand{\xxxParagraphStar}[1]{\oldparagraph*{#1}\mbox{}}
  \newcommand{\xxxParagraphNoStar}[1]{\oldparagraph{#1}\mbox{}}
\fi
\ifx\subparagraph\undefined\else
  \let\oldsubparagraph\subparagraph
  \renewcommand{\subparagraph}{
    \@ifstar
      \xxxSubParagraphStar
      \xxxSubParagraphNoStar
  }
  \newcommand{\xxxSubParagraphStar}[1]{\oldsubparagraph*{#1}\mbox{}}
  \newcommand{\xxxSubParagraphNoStar}[1]{\oldsubparagraph{#1}\mbox{}}
\fi
\makeatother


\usepackage{longtable,booktabs,array}
\usepackage{calc} % for calculating minipage widths
% Correct order of tables after \paragraph or \subparagraph
\usepackage{etoolbox}
\makeatletter
\patchcmd\longtable{\par}{\if@noskipsec\mbox{}\fi\par}{}{}
\makeatother
% Allow footnotes in longtable head/foot
\IfFileExists{footnotehyper.sty}{\usepackage{footnotehyper}}{\usepackage{footnote}}
\makesavenoteenv{longtable}
\usepackage{graphicx}
\makeatletter
\newsavebox\pandoc@box
\newcommand*\pandocbounded[1]{% scales image to fit in text height/width
  \sbox\pandoc@box{#1}%
  \Gscale@div\@tempa{\textheight}{\dimexpr\ht\pandoc@box+\dp\pandoc@box\relax}%
  \Gscale@div\@tempb{\linewidth}{\wd\pandoc@box}%
  \ifdim\@tempb\p@<\@tempa\p@\let\@tempa\@tempb\fi% select the smaller of both
  \ifdim\@tempa\p@<\p@\scalebox{\@tempa}{\usebox\pandoc@box}%
  \else\usebox{\pandoc@box}%
  \fi%
}
% Set default figure placement to htbp
\def\fps@figure{htbp}
\makeatother


% definitions for citeproc citations
\NewDocumentCommand\citeproctext{}{}
\NewDocumentCommand\citeproc{mm}{%
  \begingroup\def\citeproctext{#2}\cite{#1}\endgroup}
\makeatletter
 % allow citations to break across lines
 \let\@cite@ofmt\@firstofone
 % avoid brackets around text for \cite:
 \def\@biblabel#1{}
 \def\@cite#1#2{{#1\if@tempswa , #2\fi}}
\makeatother
\newlength{\cslhangindent}
\setlength{\cslhangindent}{1.5em}
\newlength{\csllabelwidth}
\setlength{\csllabelwidth}{3em}
\newenvironment{CSLReferences}[2] % #1 hanging-indent, #2 entry-spacing
 {\begin{list}{}{%
  \setlength{\itemindent}{0pt}
  \setlength{\leftmargin}{0pt}
  \setlength{\parsep}{0pt}
  % turn on hanging indent if param 1 is 1
  \ifodd #1
   \setlength{\leftmargin}{\cslhangindent}
   \setlength{\itemindent}{-1\cslhangindent}
  \fi
  % set entry spacing
  \setlength{\itemsep}{#2\baselineskip}}}
 {\end{list}}
\usepackage{calc}
\newcommand{\CSLBlock}[1]{\hfill\break\parbox[t]{\linewidth}{\strut\ignorespaces#1\strut}}
\newcommand{\CSLLeftMargin}[1]{\parbox[t]{\csllabelwidth}{\strut#1\strut}}
\newcommand{\CSLRightInline}[1]{\parbox[t]{\linewidth - \csllabelwidth}{\strut#1\strut}}
\newcommand{\CSLIndent}[1]{\hspace{\cslhangindent}#1}



\setlength{\emergencystretch}{3em} % prevent overfull lines

\providecommand{\tightlist}{%
  \setlength{\itemsep}{0pt}\setlength{\parskip}{0pt}}



 


\makeatletter
\@ifpackageloaded{caption}{}{\usepackage{caption}}
\AtBeginDocument{%
\ifdefined\contentsname
  \renewcommand*\contentsname{Table of contents}
\else
  \newcommand\contentsname{Table of contents}
\fi
\ifdefined\listfigurename
  \renewcommand*\listfigurename{List of Figures}
\else
  \newcommand\listfigurename{List of Figures}
\fi
\ifdefined\listtablename
  \renewcommand*\listtablename{List of Tables}
\else
  \newcommand\listtablename{List of Tables}
\fi
\ifdefined\figurename
  \renewcommand*\figurename{Figure}
\else
  \newcommand\figurename{Figure}
\fi
\ifdefined\tablename
  \renewcommand*\tablename{Table}
\else
  \newcommand\tablename{Table}
\fi
}
\@ifpackageloaded{float}{}{\usepackage{float}}
\floatstyle{ruled}
\@ifundefined{c@chapter}{\newfloat{codelisting}{h}{lop}}{\newfloat{codelisting}{h}{lop}[chapter]}
\floatname{codelisting}{Listing}
\newcommand*\listoflistings{\listof{codelisting}{List of Listings}}
\makeatother
\makeatletter
\makeatother
\makeatletter
\@ifpackageloaded{caption}{}{\usepackage{caption}}
\@ifpackageloaded{subcaption}{}{\usepackage{subcaption}}
\makeatother
\usepackage{bookmark}
\IfFileExists{xurl.sty}{\usepackage{xurl}}{} % add URL line breaks if available
\urlstyle{same}
\hypersetup{
  pdftitle={Early Identification of Human Cognitive Biases in the Design of Joint Cognitive Systems: A Feasibility Study},
  pdfauthor={Martin Schmettow; Joel Surall; Valeria Resendez Gomez},
  pdfkeywords={cognitive bias, retrieval-augmented generation, decision
support, human--AI interaction, Cognitive Task Analysis, debiasing},
  colorlinks=true,
  linkcolor={blue},
  filecolor={Maroon},
  citecolor={Blue},
  urlcolor={Blue},
  pdfcreator={LaTeX via pandoc}}


\title{Early Identification of Human Cognitive Biases in the Design of
Joint Cognitive Systems: A Feasibility Study}
\author{Martin Schmettow \and Joel Surall \and Valeria Resendez Gomez}
\date{}
\begin{document}
\maketitle
\begin{abstract}
Cognitive biases systematically impair professional decision-making, and
AI systems trained on human-generated data risk inheriting or amplifying
these distortions. This paper presents a proof-of-concept for a
literature-grounded retrieval-augmented generation (RAG) system aimed at
detecting cognitive biases in narratives collected during requirements
elicitation, with dual application goals: preventing biased patterns
from entering AI training data and enabling targeted bias mitigation in
joint human--AI operation. Across five professional domains---emergency
medicine, aviation, finance, cybersecurity, and engineering---the system
identified 20 of 21 intentionally embedded biases in structured
scenarios. Testing with simulated Cognitive Task Analysis (CTA)
interview narratives revealed key technical requirements for a RAG-based
expert system, including prompt-length limits and model capacity.
Findings indicate that retrieval-grounded approaches are a promising
starting point for AI-supported bias awareness in complex professional
contexts, contingent on curated knowledge bases and capable language
models.
\end{abstract}


\section{Introduction}\label{introduction}

\subsection{Cognitive Bias in Professional
Decision-Making}\label{cognitive-bias-in-professional-decision-making}

Professionals in domains such as healthcare, aviation, and management
frequently make critical decisions under conditions of time pressure,
uncertainty, and high workload (Berthet 2022; Egoda Kapuralalage et al.
2025; Korentsides et al. 2026). In such environments, decision-making is
often supported by heuristics---intuitive mental shortcuts---that allow
individuals to process information rapidly and efficiently (Agarwal et
al. 2023; Korteling and Toet 2022). Although heuristics help manage
limited cognitive resources, they can also lead to systematic deviations
from rational standards of reasoning and probability. These deviations
are commonly described as cognitive biases (Knipper et al. 2025;
Korteling and Toet 2022).

Cognitive biases have been linked to errors in a wide range of
professional contexts. In surgical environments they have been
implicated in a substantial proportion of adverse clinical events
(Agarwal et al. 2023). In emergency departments, biases such as
anchoring and availability can influence triage decisions and contribute
to diagnostic errors (Bhuvaneswari et al. 2025; Egoda Kapuralalage et
al. 2025). Similarly, in strategic and managerial decision-making,
confirmation bias and overconfidence can distort risk assessment and
contribute to flawed strategic choices (Rau and Bromiley 2025;
Theodorakopoulos, Theodoropoulou, and Halkiopoulos 2025). These findings
suggest that cognitive biases represent a persistent challenge for
reliable decision-making in complex professional environments.

\subsection{AI-Supported Decision-Making and Bias
Amplification}\label{ai-supported-decision-making-and-bias-amplification}

At the same time, artificial intelligence (AI) and automated
decision-support systems are increasingly integrated into professional
workflows. While such systems are often expected to improve decision
quality, they do not automatically eliminate cognitive biases. AI
systems frequently inherit patterns present in human-generated data and
can reproduce or amplify biased reasoning structures (Aich 2025; Vicente
and Matute 2023). In addition, the presence of automated recommendations
can influence human judgment through \emph{automation bias}---the
tendency to over-rely on algorithmic outputs even when contradictory
information is available (Romeo and Conti 2026). Under conditions of
time pressure and cognitive load, professionals may therefore defer
uncritically to algorithmic suggestions, creating situations in which
human and machine biases interact rather than cancel each other out (Lim
2025; Mahajan et al. 2025).

In response to these challenges, recent research has explored debiasing
strategies aimed at promoting more reflective and analytically engaged
decision-making. These approaches include metacognitive training,
cognitive forcing functions, and system designs that introduce
deliberate friction to encourage users to reconsider initial judgments
(Baruh, Cemalcilar, and Stoycheff 2026; Cabitza et al. 2024). Other
proposals involve collaborative AI architectures, such as multi-agent
systems or adversarial prompting approaches, which challenge preliminary
conclusions and encourage the consideration of alternative explanations
(Cabitza et al. 2024; Ke et al. 2024). Across these approaches, a common
principle is that decision quality can improve when individuals are
supported in questioning initial interpretations and considering
alternative explanations.

\subsection{Retrieval-Augmented Generation for Bias
Detection}\label{retrieval-augmented-generation-for-bias-detection}

Research on cognitive biases is rooted in the concept of bounded
rationality, first introduced by Herbert Simon. The
heuristics-and-biases research program developed by Kahneman and Tversky
later expanded this work by showing how human judgment systematically
departs from ideal rational models. The theme of Bounded Rationality has
been very productive in the past and present, which makes it challenging
to absorb into a user requirements analysis method. At the same time the
immense body of literature could serve very well to create an LLM-based
conversational agent using retrieval-augmented generation (RAG). In
contrast to unguided language model generation, RAG systems combine
generative models with external knowledge retrieval, allowing responses
to be grounded in a curated body of literature or domain-specific
knowledge (Albashrawi 2025; Soroush et al. 2025). This grounding
improves transparency and traceability by linking generated explanations
to identifiable sources. For topics such as cognitive bias detection,
this architecture is particularly relevant because meaningful support
requires connecting narrative descriptions of decision processes with
established theoretical concepts and evidence-based debiasing
strategies.

Detecting biases in realistic professional narratives, however, remains
difficult. Cognitive biases are rarely expressed explicitly; instead,
they are often reflected indirectly in interpretation patterns,
selective framing of information, or reliance on memorable narratives
rather than statistical evidence (Echterhoff et al. 2024; Korteling,
Paradies, and Sassen-van Meer 2023). These characteristics make subtle
biases difficult to identify through introspection or self-report.
Narrative materials such as Cognitive Task Analysis (CTA) interviews
therefore provide a challenging but ecologically relevant test case for
exploring AI-supported bias detection.

\subsection{Aim of the Study}\label{aim-of-the-study}

For the design og cognitive joint systems, two goals can be derived from
Bounded Rationality, (a) a \emph{preventive strategy} in which
identified biases are prevented from entering AI training data, and (b)
an \emph{progressive strategy} in which the system is specifically
trained to counteract biases in joint human--AI operation. The present
project developed and explored a proof-of-concept for a
literature-grounded RAG system aimed at detecting cognitive biases in
user requirement artifacts. The project investigated whether an AI
system anchored in a curated body of cognitive bias literature could
identify bias patterns in fictional professional scenarios and narrative
interview material across multiple professional domains.

\section{Method}\label{method}

This project followed an exploratory, design-oriented proof-of-concept
approach structured into five successive phases: (1) literature
exploration and corpus scoping, (2) automated corpus construction, (3)
RAG-based testing with structured professional scenarios, (4) evaluation
framework development using simulated CTA interviews, and (5)
implementation of a local RAG system to address technical constraints.
In the following sections the approach is summarized, whereas a detailed
report can be found in Surall and Schmettow (2026).

\subsection{Literature Corpus (Phases
1--2)}\label{literature-corpus-phases-12}

The first phase systematically identified literature on cognitive
biases, debiasing strategies and human--AI interaction. Searches covered
PsycINFO, PubMed/MEDLINE, Web of Science, Scopus, ACM Digital Library,
IEEE Xplore, and Google Scholar. Search string development was supported
iteratively using ChatGPT. In the second phase, a Python-based Scopus
web scraper was used to extract article metadata, producing an initial
corpus of more than 250 research papers after combined manual and
automated retrieval.

\subsection{Scenario Testing with NotebookLM (Phase
3)}\label{scenario-testing-with-notebooklm-phase-3}

To conduct an initial feasibility test, fictional professional scenarios
were constructed using ChatGPT across five domains: emergency medicine
(chest-pain triage), aviation (flight decision-making), finance (loan
approval), cybersecurity (incident response), and engineering/project
management. Each scenario embedded a predefined set of cognitive biases
reflecting decision conditions such as time pressure, incomplete
information, workload, and social influence. In total, 21 biases were
intentionally embedded across the five scenarios. The literature corpus
was loaded into Google NotebookLM---which supports up to 50
simultaneously embedded PDF sources---and the system was prompted to
identify cognitive biases, explain their contextual relevance, and
suggest debiasing actions. Detected biases were compared with the
embedded reference set; ChatGPT assisted in summarizing the overlap, and
results were manually reviewed for accuracy.

\subsection{CTA Interview Materials (Phase
4)}\label{cta-interview-materials-phase-4}

To increase task realism, two simulated CTA interview narratives were
generated using ChatGPT, with instructions to embed up to three biases
per narrative without making them explicitly identifiable. The first
interview was shorter and focused on operational decision-making; the
second was longer and included unrelated contextual content (stress
management, hobbies) to simulate realistic interview noise. An annotated
reference solution---mapping each bias to specific transcript passages
and explanations---was created to serve as a benchmark for future
human--RAG comparison.

Biases were embedded in statements describing operational reasoning, for
example: ``I considered rerouting it slightly, but last time I tried a
new path it didn't help, so I stuck with the standard route.'' This
passage illustrates a preference for established procedures and
reluctance to deviate from prior decisions, which corresponds to status
quo bias and conservatism. Other passages reflected patterns such as
reliance on past outcomes, assumptions about predictability of others'
behaviour, and optimistic expectations about system stability.

\subsection{Local RAG Implementation (Phase
5)}\label{local-rag-implementation-phase-5}

Because NotebookLM imposes a prompt input limit of approximately 250
words---preventing analysis of full interview transcripts---a local RAG
system was implemented using OpenNotebookLM (an open-source alternative)
running via the Ollama inference framework. Setup required manual
troubleshooting beyond the available documentation. Due to hardware
constraints (limited VRAM and RAM), only a lower-capacity language model
could be run locally. The system was evaluated for its ability to detect
biases in the embedded literature-grounded setting.

\section{Results}\label{results}

\subsection{Structured Scenario
Testing}\label{structured-scenario-testing}

Across the five professional scenarios, NotebookLM identified 20 of the
21 intentionally embedded biases (95.2\%). The single missed bias was
\emph{base-rate neglect} in the emergency department scenario---a
probabilistic reasoning bias not explicitly cued in the narrative. In
the remaining four scenarios, all embedded biases were detected. In
several cases the system applied slightly different labels than the
predefined categories, but the underlying concepts corresponded closely
to the intended biases.

Beyond identification, the system consistently linked detected biases to
contextual triggers within the scenarios, including time pressure,
reliance on automated systems, authority influence, prior experiences,
and organizational incentives. Suggested debiasing strategies focused on
slowing decision processes, introducing structured review moments,
applying checklists or cognitive forcing techniques, and encouraging
consideration of alternative interpretations. Detection was strongest
for explicitly cued biases such as anchoring, confirmation bias,
automation bias, authority bias, and escalation of commitment.

\subsection{CTA Interview Framework}\label{cta-interview-framework}

Phase 4 produced two simulated CTA interview narratives embedding a
range of subtle biases---including confirmation bias, status
quo/conservatism bias, optimism bias, availability bias, overconfidence,
and anchoring---alongside a structured evaluation framework. This
framework comprised a participant reporting form, a multi-evaluator
comparison checklist, and an annotated reference solution.

However, the planned comparison between RAG-based and human expert
detection could not be executed. The prompt-length limit of NotebookLM
(approximately 250 words) prevented the full interview transcripts from
being submitted as input. Fragmenting the interviews would have removed
contextual cues necessary for interpreting the embedded biases, so this
strategy was not pursued.

\subsection{Local RAG System}\label{local-rag-system}

The local OpenNotebookLM system was successfully installed and could
process document sources and generate retrieval-grounded responses.
Initial tests confirmed that the system could identify cognitive biases
in provided material. However, detection quality was lower than in the
NotebookLM experiments: some biases were described using broader
conceptual language rather than precise bias labels, and some embedded
biases were missed entirely. These limitations are attributed primarily
to the lower reasoning capacity of the model required by the available
hardware.

The experiments also highlighted an important corpus-design
consideration: retrieval quality did not improve with larger document
collections. A large heterogeneous embedding space can reduce retrieval
precision, suggesting that a smaller, carefully curated set of sources
is preferable to a broad but loosely structured corpus.

\section{Discussion}\label{discussion}

\subsection{Feasibility and Scope}\label{feasibility-and-scope}

The results demonstrate that a literature-grounded RAG system can
operationalize cognitive bias detection when narrative cues are
sufficiently salient. The near-complete detection rate in structured
scenarios (20/21 biases) supports the feasibility of the approach for
professional decision contexts. The single missed bias---base-rate
neglect---was consistent with the broader literature on probabilistic
reasoning biases, which are harder to detect because they depend on
implicit statistical reasoning rather than explicit narrative cues
(Knipper et al. 2025).

The ability to link detected biases to contextual decision factors and
to generate plausible debiasing suggestions illustrates the potential of
RAG systems to translate theoretical constructs from cognitive
psychology into interpretive tools for professional decision situations.
These properties are relevant to both application strategies motivating
the project: preventing biased patterns from contaminating AI training
corpora and supporting active debiasing in joint human--AI operation.

\subsection{Technical Constraints}\label{technical-constraints}

Several practical constraints limit the current scope of the system. The
prompt-length limitation of cloud-based RAG interfaces prevented the
analysis of full CTA interview transcripts, highlighting a critical gap
between the complexity of realistic qualitative data and the technical
affordances of current tools. The local RAG implementation partially
addressed this gap but introduced trade-offs related to hardware
capacity and model performance. Taken together, these findings suggest
that the effectiveness of literature-grounded RAG systems depends on
three interacting factors: the reasoning capabilities of the underlying
language model, system-level constraints such as context-length limits,
and the quality and structure of the embedded knowledge corpus.

\subsection{Limitations}\label{limitations}

The professional scenarios used in testing were fictional and relatively
structured, limiting ecological validity. Because cognitive biases were
intentionally embedded within the scenarios, their presence may have
been more detectable than in naturally occurring professional
narratives. As a result, the detection rates observed in structured
testing conditions may overestimate performance in real-world settings
where biases are more subtle and context-dependent. On teh opposite, the
simulated CTA interviews were not fully evaluated due to the
prompt-length constraint. The local RAG setup underestimates achievable
performance because hardware limitations required a low-capacity model.
Finally, the literature corpus was only partially curated, and the
subset available for embedding may not have been optimally composed for
targeted retrieval.

\subsection{Future Directions}\label{future-directions}

Future work should (1) develop a focused, curated corpus of key
publications on cognitive biases and decision-making to improve
retrieval precision; (2) test the system using more capable local RAG
configurations on higher-performance hardware, enabling analysis of
complete CTA interview transcripts without fragmentation; (3) conduct a
systematic comparison between human expert bias detection and RAG-based
detection using the evaluation framework developed in Phase 4; and (4)
explore integration of RAG-based bias detection into professional
training environments as a reflective decision-support tool.

In addition, future studies should evaluate system performance using
more ecologically valid materials, including real-world professional
narratives and naturally occurring CTA interview data. Such studies
should incorporate systematic reporting of false positives, false
negatives, and cases of label disagreement to provide a more
comprehensive understanding of detection behavior. Establishing
inter-rater reliability among human annotators will further strengthen
the validity of reference annotations.

Future evaluations should also include human comparison baselines,
allowing calculation of standard performance metrics such as precision
and recall. These measures will support more rigorous benchmarking and
enable clearer assessment of how well detection accuracy observed in
synthetic materials generalizes to real-world contexts. Finally, future
research may explore adaptive retrieval strategies and corpus
structuring methods to optimize retrieval performance in large-scale
knowledge bases.

\section{Conclusion}\label{conclusion}

In an ideal joint cognitive system the partners complement each other
rather than mirror or in the worst case amplify each others biases.
Being able to identify potential biases early in the process is a
critical step towards this goal. The present project demonstrates that
detecting cognitive biases in requirements narratives is possible,
making use of the extensive scientific literature. Bridging the gap
between scientific research and engineering has always been a challenge.
It seems feasible to build RAG-based tools that deliver on-demand
scientific knowledge to system analysts in the field.

\section*{References}\label{references}
\addcontentsline{toc}{section}{References}

\phantomsection\label{refs}
\begin{CSLReferences}{1}{0}
\bibitem[\citeproctext]{ref-agarwal2023}
Agarwal, S., H. Aylmore, H. Marcus, and A. Pandit. 2023. {``374
Cognitive Biases and Heuristics in Surgical Settings: {A} Systematic
Review.''} \emph{British Journal of Surgery} 110 (Supplement\_7):
znad258.011. \url{https://doi.org/10.1093/bjs/znad258.011}.

\bibitem[\citeproctext]{ref-aich2025}
Aich, P. 2025. {``Cognitive Bias in {AI} Recommendations:
{Understanding} and Mitigating Human-{AI} Decision-Making Errors.''} In
\emph{Proceedings of the 13th International Conference on Human-Agent
Interaction}, 583--84. \url{https://doi.org/10.1145/3765766.3765887}.

\bibitem[\citeproctext]{ref-albashrawi2025}
Albashrawi, M. 2025. {``Generative {AI} for Decision-Making: {A}
Multidisciplinary Perspective.''} \emph{Journal of Innovation \&
Knowledge} 10 (4): 100751.
\url{https://doi.org/10.1016/j.jik.2025.100751}.

\bibitem[\citeproctext]{ref-baruh2026}
Baruh, L., Z. Cemalcilar, and A. Stoycheff. 2026. {``Mitigating
Cognitive Biases in Higher Education Contexts Through a Brief
Video-Based Intervention: {A} Randomized Controlled Trial.''}
\emph{Computers \& Education} 245: 105535.
\url{https://doi.org/10.1016/j.compedu.2025.105535}.

\bibitem[\citeproctext]{ref-berthet2022}
Berthet, V. 2022. {``The Impact of Cognitive Biases on Professionals'
Decision-Making: {A} Review of Four Occupational Areas.''}
\emph{Frontiers in Psychology} 12: 802439.
\url{https://doi.org/10.3389/fpsyg.2021.802439}.

\bibitem[\citeproctext]{ref-bhuvaneswari2025}
Bhuvaneswari, E., K. D. V. Prasad, M. Ashraf, S. Jadhav, T. R. K. Rao,
and T. Sudha Rani. 2025. {``A Human-Centered Hybrid {AI} Framework for
Optimizing Emergency Triage in Resource-Constrained Settings.''}
\emph{Intelligence-Based Medicine} 12: 100311.
\url{https://doi.org/10.1016/j.ibmed.2025.100311}.

\bibitem[\citeproctext]{ref-cabitza2024}
Cabitza, F., C. Natali, L. Famiglini, A. Campagner, V. Caccavella, and
E. Gallazzi. 2024. {``Never Tell Me the Odds: {Investigating} Pro-Hoc
Explanations in Medical Decision Making.''} \emph{Artificial
Intelligence in Medicine} 150: 102819.
\url{https://doi.org/10.1016/j.artmed.2024.102819}.

\bibitem[\citeproctext]{ref-echterhoff2024}
Echterhoff, J. M., Y. Liu, A. Alessa, J. McAuley, and Z. He. 2024.
{``Cognitive Bias in Decision-Making with {LLMs}.''} In \emph{Findings
of the Association for Computational Linguistics: {EMNLP} 2024},
12640--53. \url{https://doi.org/10.18653/v1/2024.findings-emnlp.739}.

\bibitem[\citeproctext]{ref-egoda2025}
Egoda Kapuralalage, T. N., H. F. Chan, U. Dulleck, J. A. Hughes, B.
Torgler, and S. Whyte. 2025. {``Clinical Decision-Making: {Cognitive}
Biases and Heuristics in Triage Decisions in the Emergency
Department.''} \emph{The American Journal of Emergency Medicine} 92:
60--67. \url{https://doi.org/10.1016/j.ajem.2025.02.043}.

\bibitem[\citeproctext]{ref-ke2024}
Ke, Y., R. Yang, S. A. Lie, T. X. Y. Lim, Y. Ning, I. Li, H. R.
Abdullah, D. S. W. Ting, and N. Liu. 2024. {``Mitigating Cognitive
Biases in Clinical Decision-Making Through Multi-Agent Conversations
Using Large Language Models: {Simulation} Study.''} \emph{Journal of
Medical Internet Research} 26: e59439.
\url{https://doi.org/10.2196/59439}.

\bibitem[\citeproctext]{ref-knipper2025}
Knipper, R. A., C. S. Knipper, K. Zhang, V. Sims, C. Bowers, and S.
Karmaker. 2025. {``The Bias Is in the Details: {An} Assessment of
Cognitive Bias in {LLMs}.''}
\url{https://doi.org/10.48550/arxiv.2509.22856}.

\bibitem[\citeproctext]{ref-korentsides2026}
Korentsides, J., E. R. Merwin, L. Berger, L. Laskey, S. R. Winter, B.
Sobel, and J. R. Keebler. 2026. {``The Use of Artificial Intelligence
({AI}) in the Flight Deck: {Enhancing} Human-{AI} Teamwork in
Aviation.''} \emph{Journal of the Air Transport Research Society} 6:
100099. \url{https://doi.org/10.1016/j.jatrs.2025.100099}.

\bibitem[\citeproctext]{ref-korteling2023}
Korteling, J. E., G. L. Paradies, and J. P. Sassen-van Meer. 2023.
{``Cognitive Bias and How to Improve Sustainable Decision Making.''}
\emph{Frontiers in Psychology} 14: 1129835.
\url{https://doi.org/10.3389/fpsyg.2023.1129835}.

\bibitem[\citeproctext]{ref-korteling2022}
Korteling, J. E., and A. Toet. 2022. {``Cognitive Biases.''} In
\emph{Encyclopedia of Behavioral Neuroscience}, 2nd ed., 610--19.
Elsevier. \url{https://doi.org/10.1016/B978-0-12-809324-5.24105-9}.

\bibitem[\citeproctext]{ref-lim2025}
Lim, C. 2025. {``{DeBiasMe}: {De-biasing} Human-{AI} Interactions with
Metacognitive {AIED} Interventions.''}
\url{https://doi.org/10.48550/arXiv.2504.16770}.

\bibitem[\citeproctext]{ref-mahajan2025}
Mahajan, A., Z. Obermeyer, R. Daneshjou, J. Lester, and D. Powell. 2025.
{``Cognitive Bias in Clinical Large Language Models.''} \emph{Npj
Digital Medicine} 8 (1): 428.
\url{https://doi.org/10.1038/s41746-025-01790-0}.

\bibitem[\citeproctext]{ref-rau2025}
Rau, D., and P. Bromiley. 2025. {``A Review of Cognitive Biases in
Strategic Decision Making.''} \emph{Long Range Planning} 58 (3): 102529.
\url{https://doi.org/10.1016/j.lrp.2025.102529}.

\bibitem[\citeproctext]{ref-romeo2026}
Romeo, G., and D. Conti. 2026. {``Exploring Automation Bias in
Human--{AI} Collaboration: {A} Review and Implications for Explainable
{AI}.''} \emph{AI \& Society} 41 (1): 259--78.
\url{https://doi.org/10.1007/s00146-025-02422-7}.

\bibitem[\citeproctext]{ref-soroush2025}
Soroush, A., M. Giuffrè, S. Chung, and D. L. Shung. 2025. {``Generative
Artificial Intelligence in Clinical Medicine and Impact on
Gastroenterology.''} \emph{Gastroenterology} 169 (3): 502--517.e1.
\url{https://doi.org/10.1053/j.gastro.2025.03.038}.

\bibitem[\citeproctext]{ref-surall2026}
Surall, J., and M. Schmettow. 2026. {``Cognitive Biases in Human-{AI}
Interaction: {A} Systematic Review of Bias Types, Mitigation Strategies,
and Future Research Directions.''}
\url{https://github.com/schmettow/pub-cog-debias/blob/main/Surall_Schmettow_cognitive_debiasing.pdf}.

\bibitem[\citeproctext]{ref-theodorakopoulos2025}
Theodorakopoulos, L., A. Theodoropoulou, and C. Halkiopoulos. 2025.
{``Cognitive Bias Mitigation in Executive Decision-Making: {A}
Data-Driven Approach Integrating Big Data Analytics, {AI}, and
Explainable Systems.''} \emph{Electronics} 14 (19): 3930.
\url{https://doi.org/10.3390/electronics14193930}.

\bibitem[\citeproctext]{ref-vicente2023}
Vicente, L., and H. Matute. 2023. {``Humans Inherit Artificial
Intelligence Biases.''} \emph{Scientific Reports} 13 (1): 15737.
\url{https://doi.org/10.1038/s41598-023-42384-8}.

\end{CSLReferences}




\end{document}
